\scanpage{76}
\mspage{20}

\uline{Corollaire 1} Si $X\xrightarrow{f}Y$ de $\mathrm{Cat}_{/S}$ est dans $W$, alors $\deleted{\unsure{\Phi}}\Psi(X)\xrightarrow{\Psi(f)}\Psi(Y)$ aussi. Mais on se gardera de croire que c'est dans $W_S$ \unsure{du|au} sens de $\underline{\mathrm{Cofib}}\,S$ ou $\underline{\mathrm{Cofib}\,\mathrm{sc}}\,S$.

Car sur les fibres \inserted{de $\Psi(X),\Psi(Y)$} \deleted{c'est} on trouve les flèches
\[ f_{/s}:X_{/s}\unsure{\longrightarrow|\text{---}}Y_{/s} \]
induites par $f$, qui n'ont aucune raison d'être dans $W$. Qu'elles le soient \unsure{signifie|signifierait} que $f:X\unsure{\text{---}|\to}Y$ est \unsure{dans} $W$ ``localement au dessus de $S$'', ce qui \inserted{(pour les $W$ qui nous intéressent)} est beaucoup plus fort que d'être dans $W$. (Si $Y=S$, cela signifie en fait que $p:X\to S$ est $W$-asphérique, beaucoup plus fort que \deleted{\unsure{le fait}} $p\in W$.)
% slanted
\margin{cf \textcircled{1} \unsure{$\blacksquare$} page 15.}

Donc \deleted{$\unsure{\Psi}$} pour prouver que
\[ \unsure{\text{``}}\,\underline{\mathrm{Cofib}}\,S\to\mathrm{Cat}_{/S}\quad\text{ou}\quad \underline{\mathrm{Cofib}\,\mathrm{sc}}\,S\unsure{\text{---}|\to}\mathrm{Cat}_{/S} \]
induisent des équivalences pour les catégories localisées, le foncteur $\Psi$ directement ne \deleted{\unsure{nous}} sera pas utilisable. Cependant\,:

\uline{\unsure{Corollaire} 2} Supposons $X,Y$ dans $\underline{\mathrm{Prop}}\,S$ \inserted{ou dans $\mathrm{Fib\,hot}\,S$}, et considérons $\Psi(f):\Psi(X)\to\Psi(Y)$, pour que ce soit une $W$-équivalence dans $\underline{\mathrm{Cofib}\,\mathrm{sc}}\,S$ \deleted{(ou dans $\underline{\mathrm{Cofib}}\,S$, ce qui revient au même),}

% archive sheet number: 75
